In this paper, KDS-EL is proposed to overcome search blindness, evolutionary amnesia, destructive perturbation, and static weighting bottlenecks in ultra-high-dimensional wrapper feature selection. By integrating mutual information prior initialization with adaptive recovery, KDS-EL establishes informative search basins while safeguarding diversity. To suppress destructive perturbations, the KER strategy harvests fitness differentials under exponential memory decay, steering targeted refinement within bounded computation. Furthermore, dynamic Two-Phase Weight Scheduling temporally decouples coarse dimensional pruning from fine classification refinement. Benchmark evaluations across 20 datasets demonstrate that KDS-EL delivers competitive overall performance under multi-criteria evaluation, attaining 0.9574 mean accuracy, contracting feature subsets to an average of 87.8 dimensions, and achieving a significant 3.4-fold acceleration over the original MEL algorithm.

Beyond nearest-neighbor metrics, extensive evaluations confirm that the discovered compact feature subsets exhibit robust cross-model discriminative transferability across Support Vector Machines and Random Forest models. On high-dimensional microarrays, KDS-EL suppresses feature bloat, isolating compact candidate gene subsets of 39 to 83 dimensions that successfully enrich for established oncological markers, including \textit{COL11A1}, \textit{GUCA2A}, \textit{HPN}, and \textit{AMACR}. Methodologically, the synergy of knowledge-guided importance accumulation, bounded elite refinement, and dynamic objective scheduling provides a versatile foundational paradigm for continuous-discrete swarm intelligence.

Several promising trajectories emerge for future exploration. First, integrating self-supervised embeddings could discover higher-order epistatic interactions across massive cohorts. Second, cross-cohort transfer mechanisms could enable reusing accumulated importance memories. Third, evaluating cross-run feature selection stability alongside prospective multi-center cohort validation will be essential to advance computationally identified candidate signatures toward prospective clinical evaluation.
